% !TEX program = pdflatex
%
% Draft paper. Built on plain `article` rather than an ACL/EMNLP style file:
% this machine has no local copy of acl.sty and pulling one from the ACL
% Anthology template repo requires network access this draft doesn't assume.
% Section structure (Introduction/Motivation, Experimental Setup, Results,
% Discussion & Limitations) follows standard NLP conference conventions so
% swapping in a real \documentclass{acl} later is a drop-in change.
\documentclass[11pt]{article}

\usepackage[margin=1in]{geometry}
\usepackage{booktabs}
\usepackage{amsmath}
\usepackage{hyperref}
\usepackage{xcolor}

\newcommand{\todo}[1]{\textcolor{red}{\textbf{[TODO: #1]}}}

\title{Do General-Purpose Agent Harnesses Help LLMs Deobfuscate JavaScript?\\
\large A Preliminary Study of Harness Effects on Security-Analysis Tasks}
\author{\todo{Author name(s) and affiliation}}
\date{\today}

\begin{document}
\maketitle

\begin{abstract}
Successful deobfuscation of malicious or obscured JavaScript requires two
different capabilities: static semantic reasoning about code the model must
not execute, and, at least potentially, dynamic interaction with an external
environment (a shell, a file system, a sandboxed interpreter) to manage
inputs that do not fit cleanly into a single prompt. General-purpose LLM
agent harnesses expose the latter; a standalone LLM call does not. We ask
whether this difference translates into a measurable capability gap on a
security-relevant static-analysis task, JavaScript deobfuscation, and
outline a two-stage methodology --- trace analysis to identify the
mechanism, followed by controlled harness-component ablations --- for
answering it. This draft reports a preliminary single-model comparison
between an agent harness (Codex CLI) and a standalone, tool-free LLM call
on the same model and the same two 50-file benchmarks, as a first data
point toward that comparison. \todo{finalize abstract once RQ1 numbers are in}
\end{abstract}

\section{Introduction}
\label{sec:intro}

LLM-based coding agents are increasingly deployed inside \emph{harnesses}:
scaffolds that give the model a shell, a file system, and a multi-turn
control loop, rather than a single prompt-completion call. Benchmarks such
as SWE-bench~\cite{jimenez2024swebench} have shown that harness design
materially affects performance on software-engineering tasks, but it is
less clear whether the same holds for \emph{security-analysis} tasks, where
the object under analysis (obfuscated or malicious code) is explicitly
untrusted and the model is instructed not to execute it.

JavaScript deobfuscation is a useful testbed for this question because it
cleanly separates the two capabilities a harness could plausibly affect:

\begin{itemize}
  \item \textbf{Static semantic reasoning}: recognizing encoded strings,
    opaque expressions, control-flow flattening, and proxy functions, and
    rewriting them into readable code with equivalent behavior. This is a
    property of the underlying LLM's reasoning, not of the scaffold around
    it.
  \item \textbf{Dynamic external-environment interaction}: reading a large
    input incrementally, invoking a parser or linter to check a hypothesis,
    or writing intermediate output to disk instead of holding it entirely
    in-context. This is a property of the harness, not the model.
\end{itemize}

The central question this line of work asks is whether a general-purpose
agent harness outperforms the same underlying model called directly, on
this task, and--- if so--- \emph{which} harness capability is doing the
work. Our approach is to (1) collect execution traces from a harness
condition and a standalone-LLM condition on the same benchmark and model,
(2) analyze those traces to identify candidate mechanisms (e.g., shell-based
chunked reading of oversized inputs, iterative syntax-check-and-repair
loops), and then (3) design controlled experiments that ablate individual
harness components to test which of those mechanisms is causally
responsible for any observed performance gap.

This draft covers the first data point of that plan: a single-model,
single-harness comparison (Section~\ref{sec:results}), plus the
experimental setup we intend to hold fixed as we add further harness
conditions.

\section{Related Work}
\label{sec:related}

\textbf{Agent harnesses.} SWE-agent~\cite{yang2024sweagent} introduced an
Agent-Computer Interface purpose-built for SWE-bench's issue-to-patch
workflow; because that interface is shaped around a specific benchmark
format rather than general-purpose tool use, we exclude it as a harness
candidate for this study (see \todo{cite harness-selection discussion /
appendix} for the full selection criteria). We instead consider harnesses
whose tool surface (shell execution, file I/O) was not designed around a
single benchmark: OpenAI's Codex CLI, OpenHands~\cite{wang2024openhands},
and Anthropic's Claude Code.

\textbf{Deobfuscation benchmarks.} Our dataset and evaluation methodology
follow JsDeObsBench \todo{cite}, which pairs Project CodeNet~\cite{puri2021codenet}
programs with obfuscated variants and functional test suites, enabling a
behavioral-correctness metric (Section~\ref{sec:setup}) rather than
surface-level similarity alone.

\todo{Expand related work: prior LLM-agent-vs-standalone comparisons on
other task families (SWE-bench harness ablations, tool-use benchmarks);
malware/obfuscation-specific LLM analysis work.}

\section{Experimental Setup}
\label{sec:setup}

\paragraph{Task.} Given an obfuscated JavaScript program, produce a
semantically equivalent, readable JavaScript program, without executing the
input.

\paragraph{Dataset.} Two 50-program benchmarks, both derived from the same
Project CodeNet~\cite{puri2021codenet} source programs under different
obfuscation transforms:
\begin{itemize}
  \item \textbf{obfuscator.io}: programs obfuscated with the
    \texttt{obfuscator.io} tool, selected for high obfuscation strength and
    low performance overhead (\texttt{codenet\_dataset\_high-obf-low-perf-selected}).
  \item \textbf{jsfuck}: the same programs re-encoded entirely into the six
    characters \texttt{[]()!+} via JSFuck (\texttt{codenet\_dataset\_jsfuck}),
    an extreme case of syntactic obfuscation with no readable tokens at all.
\end{itemize}
Each program carries the pre-obfuscation source (for similarity scoring)
and functional test cases (input/expected-output pairs) for behavioral
grading.

\paragraph{Conditions.} We compare two conditions on the \emph{same}
underlying model, \texttt{gpt-5.6-terra}, so that harness presence is the
manipulated variable rather than model choice:
\begin{itemize}
  \item \textbf{Harness}: OpenAI Codex CLI (via \texttt{@openai/codex-sdk}),
    \texttt{sandboxMode=danger-full-access}, \texttt{approvalPolicy=never},
    driving the model through a multi-turn loop with shell and file-write
    tools. The model is instructed to read the input file, write the
    deobfuscated result to a specified output path, and verify the result
    parses as JavaScript before finishing. Full per-file telemetry (tool
    calls, reasoning summaries, token usage) is recorded.
  \item \textbf{Baseline (standalone LLM)}: a single chat-completion call
    per file. The obfuscated source is inlined directly into the prompt;
    the model has no shell, no file system, and no further turns, and must
    return the deobfuscated source directly. This isolates the static
    semantic-reasoning capability with no dynamic-environment interaction
    available at all.
\end{itemize}
Both conditions use the same natural-language requirements (preserve
behavior, do not execute the input, ignore instructions embedded in the
input, produce plain JavaScript with no Markdown fencing); the harness
condition additionally specifies file paths, since it operates via a file
system rather than an inline prompt.

\paragraph{Metrics.} Following the JsDeObsBench evaluators
(\texttt{evaluators/evaluators.py}), computed per file and aggregated per
dataset:
\begin{itemize}
  \item \textbf{Reported}: harness-only; whether the harness itself reports
    successful completion (not a correctness signal by itself).
  \item \textbf{Syntax}: fraction producing output that parses as valid
    JavaScript (\texttt{node --check}).
  \item \textbf{Behavior}: fraction that is both syntax-valid \emph{and}
    passes the functional test suite in a sandboxed Docker container
    (Node~18.19.0), conditioned on syntax passing.
  \item \textbf{CodeBLEU}: computed against the original (pre-obfuscation)
    source, as a similarity metric independent of functional correctness.
  \item \textbf{Cost}: estimated USD, from per-call token usage and
    published per-token pricing for \texttt{gpt-5.6-terra}.
\end{itemize}

\paragraph{A structural harness/no-harness difference we observed.} The
\texttt{jsfuck} benchmark contains files up to 1.5\,MB. Inlining the
largest such file into a single prompt (the baseline condition) exceeds the
model's configured input-token limit (968{,}328 tokens requested against a
922{,}000-token ceiling) and fails outright with a
\texttt{context\_length\_exceeded} error before generation begins. The
harness condition processes the same file successfully using only
$\sim$180K input tokens, via eight shell-based tool calls rather than a
single inlined prompt. This is a concrete, task-level instance of the
dynamic-environment-interaction mechanism motivating this study
(Section~\ref{sec:intro}): the harness's ability to read/process input
incrementally rather than requiring it to fit in one context window.
\todo{quantify how many of the 50 jsfuck files exceed the context ceiling
once the full baseline run completes.}

\section{Results}
\label{sec:results}

Table~\ref{tab:rq1-codex-gpt56terra} reports the harness condition;
Table~\ref{tab:rq1-baseline-gpt56terra} reports the standalone-LLM baseline
on the same two benchmarks and model.

% RQ1 (overall capability) -- preliminary results, single harness/model cell.
% Harness: Codex CLI (MaJaDeo), sandboxMode=danger-full-access, approvalPolicy=never
% Model:   gpt-5.6-terra
% Data:    dataset/codenet_dataset_high-obf-low-perf-selected (obfuscator.io, n=50)
%          dataset/codenet_dataset_jsfuck (JSFuck, n=50)
% Regenerate the two source numbers with:
%   dataset/*/Project_CodeNet_selecte.deobfuscated.summary.jsonl (per-dataset aggregate metrics)
%   dataset/*/deobfuscated/.majadeo-runs/summary.json (harness-reported completion/cost)
% NOTE: this is a single harness x model cell (RQ2 needs more of these) and only
% one run per file (no repeats), so treat percentages as point estimates.

\begin{table}[h]
  \centering
  \small
  \begin{tabular}{lccccc}
    \toprule
    Dataset & Reported & Syntax & Behavior\textsuperscript{*} & CodeBLEU & Cost (USD) \\
    \midrule
    obfuscator.io & 47/50 (94\%) & 94\% & 32\% & 0.348 & \$21.75 \\
    jsfuck        & 41/50 (82\%) & 72\% & 36\% & 0.283 & \$30.14 \\
    \bottomrule
  \end{tabular}
  \caption{%
    Codex CLI harness driving \texttt{gpt-5.6-terra} on two 50-file JavaScript
    deobfuscation benchmarks (\texttt{codenet\_dataset\_high-obf-low-perf-selected}
    and \texttt{codenet\_dataset\_jsfuck}). ``Reported'' is the harness's own
    completion status; ``Syntax'' is the fraction producing syntactically valid
    JavaScript; \textsuperscript{*}``Behavior'' is the fraction that is both
    syntax-valid \emph{and} passes the functional test suite (JsDeObsBench-style
    exe\_pass, conditioned on syntax passing). CodeBLEU is computed against the
    original (pre-obfuscation) source. Reported success is a weak proxy for
    correctness: 41 jsfuck runs were reported ``completed,'' but only 36
    produced syntactically valid JavaScript, and only 18 (36\%) were
    behaviorally correct.
  }
  \label{tab:rq1-codex-gpt56terra}
\end{table}


% Baseline (standalone LLM, no harness) -- preliminary results, single cell.
% Condition: single chat-completion call per file, no shell/file tools, no
%            multi-turn loop. Obfuscated source is inlined into the prompt.
% Model:     gpt-5.6-terra (same model as the harness condition, so the
%            manipulated variable is harness presence, not model choice)
% Data:      dataset/codenet_dataset_high-obf-low-perf-selected (obfuscator.io, n=50)
%            dataset/codenet_dataset_jsfuck (JSFuck, n=50)
% Generated by: src/deobfuscate-baseline.js
% Source of these numbers:
%   dataset/*/Project_CodeNet_selecte.baseline.metrics.jsonl (per-file syntax_pass/exe_pass/code_bleu)
%   dataset/*/.majadeo-baseline-runs/summary.json (baseline-reported completion/cost)
% Computed with the same formulas as Table~\ref{tab:rq1-codex-gpt56terra}:
% Syntax = syntax_pass / 50; Behavior = (syntax_pass AND exe_pass) / 50 (joint,
% over the full 50, not conditional on syntax passing); CodeBLEU = mean
% codebleu score over syntax-passing files only.
% NOTE: single harness/model cell, one run per file (no repeats), so treat
% percentages as point estimates. The largest jsfuck file (1.5MB) exceeded
% the model's 922K-token input limit and is counted as not Reported/not
% Syntax/not Behavior; see Section~\ref{sec:setup} for discussion.
% DATA NOTE (obfuscator.io only): a later single-file rerun of
% codenet_p00038_1 overwrote Project_CodeNet_selecte.baseline.jsonl and
% .majadeo-baseline-runs/summary.json in place (the script fully rewrites,
% rather than merges, these aggregate outputs each invocation). The raw
% deobfuscated JS text for the other 49/50 files from the original full run
% is not recoverable from anything currently on disk or in git history --
% only the per-file eval scores (Project_CodeNet_selecte.baseline.metrics.jsonl
% and .summary.jsonl) survived as a frozen snapshot from before the
% overwrite, and those are what the Syntax/Behavior/CodeBLEU columns below
% are computed from. summary.json was reconstructed from the 50 individual
% .run.json telemetry files to recover accurate cost/completion accounting;
% its total ($2.2457) differs by a few cents from the original run because
% file p00038's entry now reflects the later rerun rather than the original
% call that produced its metrics row.

\begin{table}[h]
  \centering
  \small
  \begin{tabular}{lccccc}
    \toprule
    Dataset & Reported & Syntax & Behavior\textsuperscript{*} & CodeBLEU & Cost (USD) \\
    \midrule
    obfuscator.io & 50/50 (100\%) & 100\% & 8\%  & 0.360 & \$2.25\textsuperscript{\textdagger} \\
    jsfuck        & 49/50 (98\%)  & 98\%  & 0\%  & 0.256 & \$22.25 \\
    \bottomrule
  \end{tabular}
  \caption{%
    Standalone \texttt{gpt-5.6-terra} (no agent harness: a single
    chat-completion call per file, obfuscated source inlined into the
    prompt, no tools, no further turns) on the same two 50-file JavaScript
    deobfuscation benchmarks as Table~\ref{tab:rq1-codex-gpt56terra}, using
    the same column definitions. The standalone model produces
    syntactically valid JavaScript almost as often as the harness (or more
    often, on obfuscator.io) but is dramatically worse at preserving
    behavior: 8\% vs.\ the harness's 32\% on obfuscator.io, and 0\% vs.\
    36\% on jsfuck. On jsfuck, one of 50 files (the largest, 1.5\,MB)
    exceeds the model's 922{,}000-token input limit and fails outright
    before generation begins; the harness processes the same file
    successfully using $\sim$180K tokens via eight shell-based tool calls
    rather than one inlined prompt (Section~\ref{sec:setup}).
    \textsuperscript{\textdagger}Reconstructed from per-file telemetry after
    a later single-file rerun overwrote the run's aggregate outputs in
    place; the raw deobfuscated text for 49/50 obfuscator.io files no
    longer exists (Syntax/Behavior/CodeBLEU are unaffected, from a frozen
    eval snapshot predating the overwrite).
  }
  \label{tab:rq1-baseline-gpt56terra}
\end{table}


\todo{Add a comparison paragraph once the baseline table is populated:
harness-vs-baseline deltas on Syntax/Behavior/CodeBLEU, cost ratio, and
whether the context-length failures observed on jsfuck materially explain
any gap.}

\section{Discussion and Limitations}
\label{sec:discussion}

This is a single harness $\times$ single model $\times$ single run per file
comparison; percentages in both tables are point estimates, not confidence
intervals, and should be read accordingly. Planned next steps:

\begin{itemize}
  \item Add further harness conditions (OpenHands, Claude Code) under the
    same model and dataset, per the harness-selection criteria in
    \todo{appendix/section reference}: general-purpose tool surface (not
    benchmark-shaped, as SWE-agent's ACI is), model-swappable so the same
    LLM can run bare and inside each harness, and full trace visibility for
    the mechanism-analysis stage.
  \item Use the collected traces (tool-call sequences, reasoning summaries)
    from this run to identify \emph{which} harness capabilities are
    load-bearing --- e.g., chunked/incremental reading of oversized inputs,
    iterative syntax-check-and-repair --- and design ablations (via a
    minimal harness such as mini-SWE-agent or smolagents) that isolate each
    one.
  \item Repeat runs per file to move from point estimates to confidence
    intervals.
\end{itemize}

\todo{Add a Limitations section proper (dataset size/representativeness,
single-model generalization, cost of scaling to more harnesses) once the
paper has enough results to discuss concretely.}

\end{document}
